% p3-06-koopman

\section{P3 §6. Koopman and Transfer-Operator Formalism}

6. Koopman and Transfer-Operator Formalism

   Figure (fig-p3-ch6-koopman). Koopman / transfer triptych --- Koopman shift / Transfer operator /
                                          Spectral gap.

  6.1 The Symbolic Expression Space
  Let X denote the space of GOLDEN BRIDGE (Trinity-Pellis Bridge) unification representations (T, c) where T $\in$ T
  and c = (c0, c1, …) $\in$ ZN is an eventually-zero sequence. Equip X with the product
  topology (discrete on T, coordinatewise on ZN).

        Remark (Agent memory analogy). The space X has a structural resemblance to
        the Trinity Hive State Architecture described in the PhD program's Appendix K
        (Agent Memory), in which symbolic state representations are maintained as layers
        with coarse (Trinity-level) and fine (Pellis-level) components. Paper 3 borrows this
        conceptual framing but restricts all claims to the symbolic-dynamics
        interpretation; no claim is made about the AI architecture.

        Definition 6.1 (Pellis shift operator). The Pellis shift sigma : X $\rightarrow$ X acts by

  sigma(T, (c0, c1, c2, …)) = (PiL(T - c0 $\varphi$0), (c1, c2, c3, …)),

  where the action on the Trinity component updates the coarse representation by
  removing the leading Pellis coefficient and rescaling the remainder, and the Pellis

  component shifts by one position. Specifically, sigma truncates the leading Pellis
  coefficient, subtracts it from the Trinity anchor, and applies PiL to keep the Trinity
  component in the monomial lattice.

  6.2 The Koopman Operator

        Definition 6.2 (Koopman operator). For an observable f : X $\rightarrow$ R in some
        function space F, the Koopman operator K : F $\rightarrow$ F is

  (K f)(x) = f(sigma(x)).

  The Koopman operator is linear regardless of whether sigma is linear, and its spectral
  properties encode global dynamical structure. This operator-theoretic linearisation of
  symbolic dynamics parallels its use in fluid mechanics and ergodic theory; see Mezić
  (2005), Nonlinear Dynamics, and Budišić, Mohr, Mezić (2012), Chaos, for the modern
  framework.

        Proposition 6.3 (Formal eigenfunctions). Eigenfunctions of K with eigenvalue
        lambda satisfy f(sigma(x)) = lambda f(x). Monomials of the form f(T, c) = Ts (for s
        $\in$ C) provide formal eigenfunctions of a multiplicative variant of K: under rescaling
        by $\varphi$, Ts ↦ ($\varphi$ T)s = $\varphi$s Ts, so $\varphi$s is the eigenvalue. The set {$\varphi$s : s $\in$ C} = C
        setminus {0}, so the formal spectrum of this multiplicative Koopman operator is
        dense in the nonzero complex numbers.

  6.3 The Transfer (Ruelle--Perron--Frobenius) Operator

        Definition 6.4 (Transfer operator). For a weight function w : X $\rightarrow$ R>0, the
        transfer operator Lw : F $\rightarrow$ F is

  (Lw f)(x) = sumy $\in$ sigma-1(x) w(y) f(y).

  This is the adjoint of the Koopman operator with respect to the weight measure. It is
  the central object of Ruelle's thermodynamic formalism and of Baladi's treatment of
  positive transfer operators. In the GOLDEN BRIDGE (Trinity-Pellis Bridge) setting, it sums over all depth-1 Pellis
  coefficient choices c $\in$ Z that are consistent with a given target x.

  Canonical weight choice. Set w(y) = $\varphi$-|c0(y)|, penalising large leading Pellis
  coefficients. Under this weight:

  (Lw f)(T, c) = sumc $\in$ Z $\varphi$-|c| f((T, c) circ c),

  where (T,c) circ c denotes prepending c to the coefficient sequence c. The sum
  converges since sumc $\in$ Z $\varphi$-|c| = (1 + $\varphi$-1)/(1 - $\varphi$-1) $\cdot$ 2 - 1 = (1 + $\varphi$-1)$\varphi$-1/(1 - $\varphi$-1) $\cdot$ …;
  more precisely, sumc $\in$ Z $\varphi$-|c| = 1 + 2sumc=1^$\in$fty $\varphi$-c = 1 + 2$\varphi$-1/(1-$\varphi$-1) = 1 +
  2psi/(1-psi) = 1 + 2psi/psi-1$\cdot$ (…). A direct calculation gives sumc=-$\in$fty$\in$fty $\varphi$-|c| =
  (1+$\varphi$-1)2/(1-$\varphi$-2) = ($\varphi$+1)2/($\varphi$2-1) = $\varphi$2/($\varphi$-1)/($\varphi$+1)$\cdot$($\varphi$+1)2 = …. The key point is that
  the sum converges because $\varphi$ > 1.

  6.4 Spectral Gaps and Approximation Quality

        Conjecture 6.5 (Spectral gap; Conjecture). There exists a Banach space B $\subset$ F
        and a weight w such that Lw : B $\rightarrow$ B has a spectral gap: the leading eigenvalue rho
        > 0 is simple and the remainder of the spectrum lies in a disk of radius r < rho.

  A spectral gap would imply that Trinity approximation quality is determined by a
  finite-depth Pellis expansion; contributions from depths D > D0($\varepsilon$) are exponentially
  suppressed by (r/rho)D. This provides the operator-theoretic justification for Conjecture
  8.1 below. The proof of this conjecture would require adapting the functional-analytic
  machinery of Baladi (2000) --- specifically the spectral gap theorems for expanding
  maps on Banach spaces of functions with bounded variation or Hölder-continuous
  functions --- to the GOLDEN BRIDGE (Trinity-Pellis Bridge) setting.

  Why a spectral gap is plausible. The Pellis shift sigma is an expanding map in the Pellis
  coordinate (each step multiplies the remainder by $\varphi$ > 1), exactly the setting where
  Baladi's positive transfer operator theory applies. The main technical obstacle is the
  non-compactness of the coefficient space ZN, which prevents direct application of the
  standard compact-operator spectral theory; an appropriate Banach space of functions
  with decay conditions would be needed.

  6.5 Analogy with Thermodynamic Formalism
  The formalism is analogous to Ruelle's thermodynamic formalism. The pressure at
  inverse temperature beta is formally

  P(beta) = limD $\rightarrow$ $\in$fty (1)/(D) log sumc $\in$ {0,1}D} e-beta sumk=0D |ck| = log
  coth(beta/2),

  achieving its supremum at beta = 0 with value P(0) = +$\in$fty, consistent with the
  unrestricted Pellis expansion having infinite entropy (unconstrained digit set). The
  critical inverse temperature is betac = log $\varphi$, at which the "energy" per step equals 1
  and the pressure P(log$\varphi$) = logcoth(log$\varphi$/2) is finite. By the Trinity anchor identity
  (star), the critical energy level ebetac = $\varphi$ satisfies $\varphi$2 + $\varphi$-2 = 3, connecting the
  thermodynamic critical point to the Trinity anchor.

  This thermodynamic analogy is conceptual: the GOLDEN BRIDGE (Trinity-Pellis Bridge) system is not a physical
  statistical-mechanical system, and no physical interpretation is asserted.

\subsection{References}

- Vasilev, D. \textit{Flos Aureus Monograph: Trinity Constants}. DOI \href{https://zenodo.org/doi/10.5281/zenodo.19227877}{10.5281/zenodo.19227877}.
- CODATA 2018 recommended values: NIST \href{https://physics.nist.gov/cuu/Constants/}{https://physics.nist.gov/cuu/Constants/}.
- Heyrovska, R. Golden-angle FSC publications --- Heyrovský Institute of Physical Chemistry.
